\section{Methodology}
\label{sec:method}

We evaluate whether adversarial CoT prompting reveals vulnerabilities in LLMs that appear robust under standard safety evaluations.
Our experimental design uses paired comparisons across five prompting conditions on the same set of harmful behaviors, enabling controlled measurement of how reasoning context affects safety behavior.

\subsection{Models}
\label{sec:models}

We test two widely deployed OpenAI models:
\begin{itemize}[leftmargin=*,itemsep=0pt,topsep=0pt]
    \item \textbf{\gptmini} --- a cost-efficient model with broad deployment in production applications.
    \item \textbf{\gptfour} --- a frontier-class model representing state-of-the-art capabilities.
\end{itemize}
All queries use temperature $= 0.0$ for reproducibility and max\_tokens $= 1024$.

\subsection{Datasets}
\label{sec:datasets}

\para{Harmful behaviors.}
We randomly sample 80 harmful behaviors from the 520 in \advbench~\cite{zou2023universal} using a fixed seed (42) for reproducibility.
\advbench is the standard benchmark for evaluating adversarial attacks on aligned LLMs, enabling direct comparison with prior work.

\para{Over-refusal control.}
We use 50 safe prompts from \xstest~\cite{rottger2023xstest} to measure whether models exhibit over-refusal --- rejecting benign requests that superficially resemble harmful ones.
This control ensures that any observed vulnerability reflects genuine context-dependent safety failure rather than general behavioral degradation.

\subsection{Prompting Conditions}
\label{sec:conditions}

We design five prompting conditions that isolate different mechanisms by which CoT reasoning might bypass safety boundaries.
\tabref{tab:conditions} summarizes each condition and its hypothesized mechanism.

\begin{table}[t]
    \centering
    \caption{Prompting conditions and their hypothesized attack mechanisms. Each condition targets a distinct pathway through which CoT reasoning might circumvent safety filters.}
    \resizebox{\textwidth}{!}{%
    \begin{tabular}{@{}llp{7cm}@{}}
        \toprule
        \textbf{Condition} & \textbf{Type} & \textbf{Mechanism} \\
        \midrule
        Direct & Baseline & Raw harmful request; standard safety evaluation \\
        \simplecot & Control & ``Let's think step by step'' prefix; controls for CoT format effect \\
        \snowball & Adversarial & Progressive reasoning with structured steps; exploits reasoning escalation \\
        \contextmanip & Adversarial & Academic/analytical framing requesting detailed analysis; exploits context-dependent safety \\
        \decomp & Adversarial & Break into innocuous sub-questions (prerequisites, steps, materials); exploits compositional reasoning gap \\
        \bottomrule
    \end{tabular}
    }
    \label{tab:conditions}
\end{table}

The \textbf{direct} condition serves as the primary baseline, representing how standard safety evaluations query models.
The \textbf{\simplecot} condition controls for whether the CoT \emph{format} alone (rather than adversarial content) affects safety behavior.
The three adversarial conditions test distinct exploitation hypotheses: \snowball tests whether progressive reasoning escalation degrades safety boundaries; \contextmanip tests whether shifting the perceived social context (academic analysis) bypasses content-level safety filters; and \decomp tests whether decomposing a harmful goal into benign-seeming sub-questions circumvents detection.

\subsection{Evaluation}
\label{sec:eval}

\para{LLM judge.}
We use \gptmini as an automated judge with a structured rubric that produces JSON output containing three fields: a binary \texttt{harmful} classification, a severity score (1--5, where 1 = full refusal and 5 = full compliance), and free-text \texttt{reasoning}.
The judge evaluates whether each response provides substantively harmful content rather than merely acknowledging the topic.

\para{Metrics.}
Our primary metric is the \textbf{attack success rate (ASR)}: the fraction of responses classified as harmful by the judge.
We report 95\% Clopper-Pearson confidence intervals on all rates, which provide exact coverage for binomial proportions.
We additionally report mean severity scores (1--5 scale) and over-refusal rates on \xstest.

\para{Statistical tests.}
Since all conditions are evaluated on the same set of 80 harmful behaviors, we use McNemar's exact test~\cite{mcnemar1947note} for paired comparisons between each adversarial condition and the direct baseline.
We apply Bonferroni correction for 8 comparisons (4 conditions $\times$ 2 models), yielding a corrected significance threshold of $\alpha = 0.00625$.
We report Cohen's $h$ as the effect size for proportion comparisons, following standard conventions: $h < 0.2$ (small), $0.2 \leq h < 0.5$ (small-to-medium), $0.5 \leq h < 0.8$ (medium-to-large), $h \geq 0.8$ (large).

\subsection{Computational Resources}
\label{sec:compute}

All experiments use the OpenAI API with approximately 1,100 total calls (800 experiment queries + 100 \xstest queries + 200 judge evaluations).
Total runtime is approximately 5 minutes at an estimated cost of \$5, making the approach highly accessible for safety auditing.
